\documentclass{article}
\usepackage{amsmath, amsthm, amsfonts, amssymb}
\theoremstyle{plain}
\newtheorem{thm}{Theorem}
\numberwithin{thm}{section}

\theoremstyle{plain}
\newtheorem{prop}{Proposition}
\numberwithin{prop}{section}

\theoremstyle{definition}
\newtheorem{defn}{Definition}
\numberwithin{defn}{section}

\theoremstyle{remark}
\newtheorem*{rem}{Remark}

\newtheorem*{cor}{Corollary}

\numberwithin{equation}{section}

\newcommand{\op}{\prime}
\begin{document}
\title{Sequences and Series}
\author{Amey Joshi}
\maketitle
\section{Sequences of real numbers}\label{s1}
\begin{defn}\label{s1d1}
A sequence is a mapping $\mathbb{N} \mapsto \mathbb{R}$.
\end{defn}

\begin{defn}\label{s1d2}
A sequence $\{x_n\}$ is said to converge to a limit $L$ if for a given 
$\epsilon > 0$, we can find $N \in \mathbb{N}$ such that for all $n > N$, 
$|x_n - L| < \epsilon$.
\end{defn}

\begin{defn}\label{s1d3}
A sequence $\{x_n\}$ is said to be bounded from above (below) if there exists
$M$ ($m$) such that $x_n \le M$ ($x_n \ge m$) for all $n \in N$.
\end{defn}

\begin{defn}\label{s1d4}
A sequence $\{x_n\}$ is said to be bounded if it is bounded from above and below.
\end{defn}

\begin{defn}\label{s1d5}
A sequence $\{x_n\}$ is said to be monotone increasing if $x_n \le x_{n+1}$ for
all $n \in \mathbb{N}$. It is said to be strictly monotone increasing if $x_n <
x_{n+1}$ for all $n \in \mathbb{N}$.
\end{defn}

\begin{defn}\label{s1d6}
A sequence $\{x_n\}$ is said to be monotone decreasing if $x_n \ge x_{n+1}$ for
all $n \in \mathbb{N}$. It is said to be strictly monotone decreasing if $x_n >
x_{n+1}$ for all $n \in \mathbb{N}$.
\end{defn}

\begin{defn}\label{s1d7}
A sequence $\{x_n\}$ is said to be a Cauchy sequence if for a given $\epsilon
> 0$, we can find $N \in \mathbb{N}$ such that for all $m, n > N$, $|x_m - x_n| <
\epsilon$.
\end{defn}

\begin{defn}\label{s1d8}
If $\{n_1, n_2, \ldots\}$ is a subset of natural numbers such that $n_1 < n_2
< \ldots$ then $\{x_{n_1}, x_{n_2}, \ldots\}$ is said to be a subsequence of 
$\{x_n\}$.
\end{defn}

\begin{defn}\label{s1d9}
If $\{x_n\}$ is a sequence of reals then
\[
\limsup x_n = \inf_{n \ge 1}(\sup_{m \ge n} x_m)
\]
and
\[
\liminf x_n = \sup_{n \ge 1}(\inf_{m \ge n} x_m)
\]
\end{defn}

\begin{prop}\label{s1p1}
If $\{x_n\}$ is convergent then it has a unique limit.
\end{prop}
\begin{proof}
Let if possible, $\lim x_n = L$ and $\lim x_n = M$. Then for a given $\epsilon
> 0$, we can find, $N_1, N_2 \in \mathbb{N}$ such that for all $n > N_1$, 
$|x_n - L| < \epsilon/2$ and for all $n > N_2$, $|x_n - M| < \epsilon/2$. Then, 
if $N = \max(N_1, N_2)$, $|L - M| = |L - x_n + x_n - M| \le |L - x_n| + 
|M - x_n| < \epsilon$.
\end{proof}

\begin{thm}[Monotone Convergence Theorem]\label{s1t1}
If $\{x_n\}$ is a monotone increasing (decreasing) sequence that is bounded
above (below) then $\{x_n\}$ is convergent and its limit is its supremum 
(infimum).
\end{thm}
\begin{proof}
If $\{x_n\}$ is a monotone increasing functions that is bounded above then the
set of the sequence values is a set of real numbers that is bounded above. By 
the completeness property of real numbers, it has a supremum $M \in \mathbb{R}$.
Fix an $\epsilon > 0$. Then there is at least one $x_m$ such that $M - \epsilon
< x_m$. For if there were none such then $M$ would not be the supremum. Since
our sequence is monotone increasing, this means that for all $n > m$, $|M - x_n|
< \epsilon$ so that $M = \lim x_n$.

The proof for monotone decreasing sequence is similar.
\end{proof}

\begin{thm}[Bolzano-Weierstrass Theorem]\label{s1t2}
Every bounded sequence has a convergent subsequence.
\end{thm}
\begin{proof}
Let $m$ and $M$ be such that $m \le x_n \le M$ for all $n$ in a sequence 
$\{x_n\}$. An index $n$ is called a peak if $x_n > x_{n + i}$ for all $i \in
\mathbb{N}$.If there are an infinite number of peaks $n_1 < n_2 < \ldots$ then
$x_{n_1} > x_{n_2} > x_{n_3} \ldots$. If there are an infinite number of peaks 
then $\{x_{n_1}, x_{n_2}, x_{n_3}, \ldots$ forms a monotone decreasing sequence
that is bounded from below. By monotone convergence theorem, \ref{s1t1}, it 
converges to its infimum. 

If the process terminates at say $n_k = N$ or if there are no peaks, in which 
casei $N=0$, pick any $n_1 > N$. Since $n_1$ is not a peak, there exists $n_2 
> n_1$ such that $x_{n_2} > x_{n_1}$. Since $n_2$ is not a peak, we can find 
$n_3 > n_2$ such that $x_{n_3} > x_{n_2}$. This gives us a monotone increasing
sequence that is bounded from above.(This process will not terminate because
we are already past the last peak in the sequence, if there was one.) By the 
monotone convergence theorem, it converges to its supremum.
\end{proof}

\begin{prop}\label{s1p2}
Every convergent sequence is a Cauchy sequence.
\end{prop}
\begin{proof}
Let $\{x_n\}$ be such that $\lim x_n = L$. Then for a fixed $\epsilon > 0$, we
can find $N \in \mathbb{N}$ such that for all $n > N$, $|x_n - L| < \epsilon/2$.
Therefore, if $m, n > N$, $|x_m - x_n| = |x_m - L + L - x_n| \le |x_m - L| + 
|x_n - L| < \epsilon$, making $\{x_n\}$ a Cauchy sequence.
\end{proof}

\begin{prop}\label{s1p3}
Every Cauchy sequence is bounded.
\end{prop}
\begin{proof}
Let $\{x_n\}$ be a Cauchy sequence. Fix an $\epsilon > 0$. There exists $N \in
\mathbb{N}$ such that for all $m, n \ge N$, $|x_m - x_n| < \epsilon$. In 
particular, for all $m > N$, $|x_m - x_N| < \epsilon$ so that for the tail of
the sequence beyond $x_N$, $x_N - \epsilon$ is the lower bound and $x_N +
\epsilon$ is the upper bound. For the finite set $\{x_1, \ldots, x_{N-1}\}$, let
$m$ and $M$ be the minimum and the maximum. Then $\min(m, x_N - \epsilon)$ is
a lower bound of the entire sequence and $\max(M, x_N + \epsilon)$ is its upper
bound, making $\{x_n\}$ a bounded sequence.
\end{proof}

\begin{prop}\label{s1p4}
Every Cauchy sequence is convergent.
\end{prop}
\begin{proof}
By proposition \ref{s1p3} a Cauchy sequence is bounded. Using Bolzano-Weierstrass
theorem, \ref{s1t2}, it has a convergent subsequence $\{x_{n_k}\}$. Let $L$ be
the limit of this subsequence. For a given $\epsilon > 0$, we can find $N_1 \in
\mathbb{N}$ such that for all $n_k > N_1$, $|x_{n_k} - L| < \epsilon/2$. For the
same $\epsilon$, since $\{x_n\}$ is a Cauchy sequence, there exists $N_2 \in
\mathbb{N}$ such that for all $m, n > N_2$, $|x_m - x_n| < \epsilon/2$. Let
$N = \max(N_1, N_2)$. Then for $m > N$ and $k$ such that $n_k > N$, $|x_m - L| 
= |x_m - x_{n_k} + x_{n_k} - L| \le |x_m - x_{n_k}| + |x_{n_k} - L| < \epsilon$.
In the last inequality $|x_m - x_{n_k}| < \epsilon/2$ by the Cauchy property
and $|x_{n_k} - L| < \epsilon/2$ because $L$ is the limit of the convergent
subsequence.
\end{proof}

\section{Tests of convergence}\label{s2}
\begin{prop}[Comparison Test]\label{s2p1}
Let $\{x_n\}$ and $\{y_n\}$ be sequences such that $x_n \le y_n$ for all $n$. 
Then:
\begin{enumerate}
\item If both sequences converge then $\lim x_n \le lim y_n$.
\item If $a_n \rightarrow \infty$ then so does $b_n$.
\item If $b_n \rightarrow -\infty$ then so does $a_n$.
\end{enumerate}
\end{prop}
\begin{proof}
Let $\lim x_n = X$ and $\lim y_n = Y$. Let, if possible, $X > Y$. Choose 
\[
\epsilon = \frac{X - Y}{2}.
\]
We can find $N_1 \in \mathbb{N}$ such that for all $n > N_1$, $|x_n - X| <
\epsilon$, that is, $X - \epsilon < x_n < X + \epsilon$. For our choice of 
$\epsilon$, this gives,
\[
x_n > \frac{X + Y}{2}.
\]
Likewise, we can find $N_2 \in \mathbb{N}$ such that for all $n > N_2$, $|y_n -
Y| < \epsilon$ or
\[
y_n < \frac{X + Y}{2}.
\]
If $N = \max(N_1, N_2)$, we have
\[
y_{N+1} < \frac{X + Y}{2} < x_{N+1}
\]
contradicting the assumption $x_n \le y_n$ for all $n$.

Now suppose $\lim x_n = \infty$. Then for any $M > 0$, we can get $N \in 
\mathbb{N}$ such that $x_N > M$. For this choice of $N$, we also have $y_N > M$. 
Since this is true for any $M$, $\lim y_n = \infty$

If $\lim y_n = -\infty$. Then for any $m < 0$, we can find $N \in \mathbb{N}$
such that $y_N < m$. For this choice of $N$, we also have $x_N < m$. Since this
is true for any negative $m$, $\lim x_n = -\infty$.
\end{proof}

\begin{prop}[The Ratio Test]\label{s2p2}
If $\{a_n\}$ is a sequence such that 
\[
\lim_{n \rightarrow \infty}\frac{a_{n+1}}{a_n} = L
\]
when $\{a_n\}$ converges if $L < 1$, diverges if $L > 1$ and the test is 
inconclusive if $L = 1$.
\end{prop}
\begin{proof}
Given an $\epsilon > 0$, we can find $N \in \mathbb{N}$ such that for $n > 
N$, 
\[
\Big|\frac{a_{n+1}}{a_n} - L\Big| < \epsilon \Rightarrow a_{n+1} < 
a_n(L + \epsilon).
\]
In particular, $a_{N + k} < a_{N}(L + \epsilon)^k$ for $k \ge 0$. Choose
$\epsilon = (1 - L)/2$ so that $L + \epsilon = (1 + L)/2 < 1$. As $k \rightarrow
\infty$, $a_{N + k} \rightarrow 0$.

If $L > 1$, the limit condition gives $a_{n+1} > a_n(L - \epsilon)$. Choose 
$\epsilon = (L - 1)/2$ so that $a_{n + 1} > a_n(L + 1)/2$. For $L > 1$, $(L+1)/2
> 1$. Furthermore, we have 
\[
a_{N+k} > a_N\left(\frac{1 + L}{2}\right)^k.
\]
As $k \rightarrow \infty$, $a_{N+k} \rightarrow \infty$.

The test is inconclusive for $L = 1$ because of:
\begin{enumerate}
\item The sequence $\{1/n\}$ converges to $0$.
\item The sequence $\{n^p\}$ diverges for any $p > 0$.
\end{enumerate}
\end{proof}

\section{Series}\label{s3}
\begin{defn}\label{s3d1}
Given a sequence $\{x_n\}$, the series $\sum x_n$ is a sequence of the partial
sums $S_n = x_1 + \cdots + x_n$. $\sum x_n$ converges or diverges according as
the sequence $\{S_n\}$ converges or diverges.
\end{defn}

\begin{defn}\label{s3d2}
A series $\sum x_n$ converges absolutely if $\sum |x_n|$ converges.
\end{defn}

\begin{prop}\label{s3p1}
An absolutely convergent series converges.
\end{prop}
\begin{proof}
Let $A_n = |x_1| + \cdots + |x_n|$ be partial sums of the series $\sum |x_n|$.
Since $\{A_n\}$ converges, for a given $\epsilon > 0$, we can find $N \in 
\mathbb{N}$ such that for $n > m > N$, $|A_n - A_m| < \epsilon \rightarrow
||x_n| + |x_{n-1}| + \cdots + |x_{m+1}|| < \epsilon \Rightarrow |x_n| + 
|x_{n-1}| + \cdots + |x_{m+1}| < \epsilon < |x_n + x_{n-1} + \cdots + x_{m+1}|
< \epsilon \Rightarrow |S_n - S_m| < \epsilon$ so that $\{S_n\}$ is a Cauchy
sequence, where $S_n$ is the partial sum of $\sum x_n$.
\end{proof}

\begin{prop}\label{s3p2}
If $\sum x_n$ converges them $\lim x_n = 0$.
\end{prop}
\begin{proof}
If $S_n = x_1 + \cdots + x_n$, convergence of $\{S_n\}$ implies that for any
$\epsilon > 0$, we can find $n > m > N$ such that $|S_n - S_m| < \epsilon$.
Choose $n = m + 1$ so that $|x_{m+1}| < \epsilon$.
\end{proof}

The contrapositive of this proposition is
\begin{prop}\label{s3p3}
If $\lim x_n = L \ne 0$, $\sum x_n$ diverges.
\end{prop}

We will now show that
\begin{prop}[Geometric series]\label{s3p4}
The geometric series $\sum r^n$ converges if $|r| < 1$ and diverges 
otherwise.
\end{prop}
\begin{proof}
If $|r| > 1, \lim |r|^n \ne 0$ so that $\sum r^n$ diverges. If $0 < r < 1$, 
\[
S_n = \frac{1 - r^n}{1 - r}
\]
and $\lim S_n = 1/(1 - r)$. If $-1 < r < 0$, $\sum r^n$ converges because 
$\sum |r|^n$ converges.
\end{proof}

\section{Tests of convergence}\label{s4}
\begin{prop}[The comparison test]\label{s4p1}
Let $\sum a_n$ and $\sum b_n$ be series of positive terms. If $a_n \le b_n$,
then $\sum a_n$ converges if $\sum b_n$ converges and $\sum b_n$ diverges if
$\sum a_n$ diverges.
\end{prop}
\begin{proof}
If $\sum b_n$ converges then for a given $\epsilon > 0$, we can find $N \in
\mathbb{N}$ such that for all $n > m > N$, $|b_n + b_{n-1} + \cdots + b_{m+1}| 
< \epsilon$. Since $0 \le a_n \le b_n$, $|a_n + a_{n-1} + \cdots + a_{m+1}| \le
|b_n + b_{n-1} + \cdots + b_{m+1}| < \epsilon$. That is $\sum a_n$ converges.

If $\sum a_n$ diverges, then for any $M > 0$ we can find $N \in \mathbb{N}$ 
such that $a_1 + \cdots + a_N > M$. Since $a_n \le b_n$, this also means that
$b_1 + \cdots + b_N \ge a_1 + \cdots + a_N > M$. Since this is true for any $M
> 0$, we see that $\sum b_n$ diverges.
\end{proof}

\begin{prop}[The $p$-series]\label{s4p2}
The $p$-series $\sum n^{-p}$ converges if $p > 1$ and diverges of $p \le 1$.
\end{prop}
\begin{proof} Consider the case $p > 1$.
\begin{eqnarray*}
\sum\frac{1}{n^p} &=& 1 + \left(\frac{1}{2^p} + \frac{1}{3^p}\right) + \\
 & & \left(\frac{1}{4^p} + \cdots + \frac{1}{7^p}\right) +
     \left(\frac{1}{8^p} + \cdots + \frac{1}{15^p}\right) + \ldots \\
 &<& 1 + \frac{1}{2^{p-1}} + \frac{1}{4^{p-1}} + \frac{1}{8^{p-1}} + \ldots
\end{eqnarray*}
This is a geometric series with first term $r = 1/2^{p-1} < 1$ whose limit is
$2^{p-1}/(2^{p-1} - 1)$. Thus $\sum n^{-p}$ converges by the comparison test.

Now consider the case $p = 1$.
\begin{eqnarray*}
\sum\frac{1}{n^p}&=&1 + \frac{1}{2} + \left(\frac{1}{3} + \frac{1}{4}\right)+\\
 & & \left(\frac{1}{5} + \cdots + \frac{1}{8}\right) +
     \left(\frac{1}{9} + \cdots + \frac{1}{16}\right) + \ldots \\
 &>& 1 + \frac{1}{2} + \frac{1}{2} + \frac{1}{2} + \frac{1}{2} + \ldots
\end{eqnarray*}
Clearly, the series $\sum n^{-1}$ diverges. If $p < 1$, $n^{-p} > n^{-1}$ so 
that by the comparison test we can conclude that $\sum n^{-p}$ diverges.
\end{proof}

\begin{prop}[The Limit Comparison Test]\label{s4p3}
Let $\sum a_n$ and $\sum b_n$ be two series of positive terms. If
\[
\lim_{n \rightarrow \infty} \frac{a_n}{b_n} = L
\]
then
\begin{enumerate}
\item If $0 < L < \infty$, either both series converge or both diverge.
\item If $L = 0$ and $\sum b_n$ converges then so does $\sum a_n$.
\item If $L = \infty$ and $\sum b_n$ diverges then so does $\sum a_n$.
\end{enumerate}
\end{prop}
\begin{proof}
We will first consider the case $0 < L < \infty$. For a given $\epsilon > 0$,
we can find $N \in \mathbb{N}$ such that for $n > N$,
\[
\Big|\frac{a_n}{b_n} - L\Big| < \epsilon \Rightarrow
b_n(L - \epsilon) < a_n < b_n(L + \epsilon).
\]
For $n > m > N$, we also have 
\[
(b_n + \cdots + b_{m+1})(L - \epsilon) < (a_n + \cdots + a_{m+1}) <
(b_n + \cdots + b_{m+1})(L + \epsilon)
\]
If $\sum b_n$ converges, for a given $\epsilon_1(L + \epsilon)$, we can find
$N_1 \in \mathbb{N}$ such that for all $n > m > N_1$, $b_n + \cdots + b_{m+1}
< \epsilon_1/(L + \epsilon)$ or that $a_n + \cdots + a_{m+1} < \epsilon_1$.
Thus if $\sum b_n$ converges then so does $\sum a_n$.

If $\sum a_n$ converges then for a given $\epsilon_1/(L - \epsilon)$ we can find
$N_2 \in \mathbb{N}$ such that for all $n > m > N_2$, $a_n + \cdots + a_{m+1}
< \epsilon_1/(L - \epsilon)$ or that $b_n + \cdots + b_{m+1} < \epsilon_1$.
Thus if $\sum a_n$ converges then so does $\sum b_n$. Since $L > 0$, we can 
always choose $\epsilon$ such that $L - \epsilon > 0$ and $\epsilon_1/(L - 
\epsilon) > 0$.

Now suppose $\sum b_n$ diverges. Then for any $M/(L - \epsilon) > 0$, we can
find $N \in \mathbb{N}$ such that for all $n>m > N$, $(b_n + \cdots + b_{m+1} > 
M/(L - \epsilon)$ or $a_n + \cdots + a_{m+1} > M$. Since $M$ is arbitrary, this 
means that $\sum a_n$ diverges.

If $\sum a_n$ diverges then for any $M(L + \epsilon) > 0$, we can find $N \in
\mathbb{N}$ such that for all $n > m >N$, $a_n + \cdots + a_{m+1} > M(L + 
\epsilon)$ or $b_n + \cdots + b_{m+1} > M$. Since $M$ is arbitrary, this means 
that $\sum b_n$ diverges.

If $L = 0$, for $\epsilon = 1$, we can find $N \in \mathbb{N}$ such that
for all $n > N$, $a_n < b_n$. Therefore, $\sum_{n > N} a_n < \sum_{n > N} b_n$.
By comparison test, $\sum_{n > N} a_n$ converges. Since $a_1 + \cdots + a_N$
is finite, $\sum a_n$ converges.

If $L = \infty$, for any $M > 1$, we can find $N \in \mathbb{N}$ such that
$a_n/b_n > M$ for all $n > N$. That is $a_n > b_nM$ for all $n > N$. Then, for
$n > m > N$, $a_n + \cdots + a_{m+1} > (b_n + \cdots + b_{m+1})M$. Since $\sum
b_n$ diverges, the left hand side of the inequality cannot be made as small
as we wish. As a result, neither can the right hand side. Therefore, $\sum a_n$
diverges.
\end{proof}

\begin{prop}[The Ratio Test]\label{s4p4}
Let $\sum a_n$ be a series of positive terms. If
\[
\lim_{n \rightarrow \infty}\frac{a_{n+1}}{a_n} = L.
\]
Then 
\begin{enumerate}
\item If $L < 1$, $\sum a_n$ converges.
\item If $L > 1$, $\sum a_n$ diverges.
\item If $L = 1$, the test is inconclusive.
\end{enumerate}
\end{prop}
\begin{proof}
If $L < 1$, we can find $r$ such that $L < r < 1$. For a given $\epsilon > 0$,
such that $\epsilon < r - L$, we can find $N \in \mathbb{N}$ such that for all 
$n > N$,
\[
\Big|\frac{a_{n+1}}{a_n} - L\Big| < \epsilon \Rightarrow a_{n+1} < 
a_n(L + \epsilon) < a_n r
\]
For $n > N$, $\sum_{n > N} a_n < \sum_{n > N} r_n$ with $r < 1$. By comparison 
test, $\sum_{n > N} a_n$ converges. Since $\sum a_n = (a_1 + \cdots + a_N) + 
\sum_{n > N} a_n$, $\sum a_n$ converges.

If $L > 1$, then for a given $\epsilon > 0$, we can find $N \in \mathbb{N}$ such
that for all $n > N$, $a_{n + 1} > a_n(L - \epsilon)$. If we choose $\epsilon$
such that $L - \epsilon > 1$, then for all $k \in \mathbb{N}$, $a_{n+k} > a_n
(L - \epsilon)^k$ for all $n > N$. Thus, the limit of $a_n$ as $n \rightarrow
\infty$ does not vanish, as a result $\sum a_n$ diverges.

If $a_n = n^{-p}$, where $p \ge 0$, 
\[
\lim_{n \rightarrow \infty}\frac{a_{n+1}}{a_n} = \lim_{n \rightarrow \infty}
\frac{n^p}{(n + 1)^p} = \lim_{n \rightarrow \infty} \frac{1}{(1 + 1/n)^p} = 1
\]
But we know that the $p$-series converges for $p > 1$ and diverges otherwise.
This illustrates the inconclusivity of the ratio test if $L = 1$.
\end{proof}

\begin{prop}[The Root Test]\label{s4p5}
Let $\sum a_n$ be a series of positive terms and $L = \limsup \sqrt[n]{a_n}$.
Then $\sum a_n$ converges if $L < 1$, diverges if $L > 1$ and is inconslusive
if $L = 1$.
\end{prop}
\begin{proof}
By the definition of $\limsup$, $L = \inf_{k \ge 1}\sup_{n > k}a_n^{1/n}$. If
$L < 1$, choose $\epsilon > 0$ and $r \in (L + \epsilon, 1)$. We can then find
$N \in \mathbb{N}$ such that for all $k > N$, $|\sup_{n>k}a_n^{1/n} - L| <
\epsilon \Rightarrow \sup_{n>k}a_n^{1/n} < L + \epsilon < r < 1$. Therefore,
for all $n > k > N$, $a_n^{1/n} < r < 1 \Rightarrow a_n < r^n < 1$. Thus, for 
all $k > N$, $a_k < r^k$ so that by comparison test, $\sum_{k > N}a_k$ 
converges. Since $a_1 + \cdots + a_N$ is finite, $\sum_{k \ge 1} a_k$ converges.

If $L > 1$, choose $\epsilon > 0$ and $r \in (1, L + \epsilon)$. We can then
find $N \in \mathbb{N}$ such that for all $k > N$, $|\sup_{n > k} a_n^{1/n} -
L| < \epsilon \rightarrow 1 < r < L + \epsilon < \sup_{n>k} a_n^{1/n}$. 
Therefore, $1 < r < a_n^{1/n}$ is true for infinitely many $n > k > N$, or
$1 < r^n < a_n$ is true for infinitely many $n > k > N$. This means that
$\lim a_n \ne 0$ and hence $\sum a_n$ diverges.

Once again consider the $p$-series with $a_n = n^{-p}$ so that $a_n^{1/n} = 
n^{-p/n} = (n^{1/n})^{-p}$ so that $\limsup\sqrt[n]{a_n} = 1$ irrespective of
$p$ because $\lim\sqrt[n]{n} = 1$. Thus, the root test is inconclusive if $L=1$.
\end{proof}

\begin{prop}[The Integral Test]\label{s4p6}
Let $f:[1, \infty) \rightarrow \mathbb{R}^+$ be a decreasing, postive-valued,
continuous function of $x$ and $a_n = f(n)$. The $\sum a_n$ converges if the 
integral
\[
\int_1^\infty f(x)dx < \infty.
\]
\end{prop}
\begin{proof}
Since $f$ is continuous, the integral can be constructed as an improper Riemann
integral. If the mesh consists of the intervals $[n, n+1]$ for all $n \ge 1$,
then
\[
a_2 + a_3 + \cdots < \int_1^\infty f(x)dx < a_1 + a_2 + a_3 + \cdots
\]
If the integral is finite, $\sum_{n \ge 2}a_n$ is finite and hence $a_1+\sum_{n
\ge 2} a_n = \sum_{n \ge 1}a_n$ is finite. That is $\sum a_n$ converges.
\end{proof}

\begin{prop}[The Alternating Series Test]\label{s4p7}
The alternating series $\sum (-1)^{n+1}a_n$, where $a_n$ are all positive,
$a_{n+1} < a_n$ for all $n$ and $\lim a_n \rightarrow 0$, converges.
\end{prop}
\begin{proof}
Let $S_n$ denote the $n$th partial sum. Then $S_{2m+2} = S_{2m} + a_{2m+1} -
a_{2m+2} > S_{2m}$ as $a_n$ are monotone decreasing positive terms. Thus, the
subsequence of even terms is monotone increasing. Likewise, $S_{2m+3} = S_{2m+1}
- a_{2m+2} + a_{2m+3} < S_{2m+1}$ so that the subsequence of odd terms is
monotone decreasing. We also have $S_{2m+1} = S_{2m} + a_{2m+1} > S_{2m}$. Thus,
\[
a_1 - a_2 \le S_{2m} < S_{2m + 1} < a_1.
\]
Thus, the subsequence of even and odd terms are monotone and bounded. Therefore,
they converge. Let $L_e$ and $L_o$ be their limits. Since $|L_o - L_e| = 
\lim_{k \rightarrow \infty} a_{2k+1} = 0$. Thus, the two limits are identical.
\end{proof}

\section{Sequences and series of functions}\label{s5}
In this section, we let $f_n$ be real-valued functions defined on a set $A$ for
every $n \in \mathbb{N}$.

\begin{defn}\label{s5d1}
The sequence $\{f_n\}$ converges to a function $f: A \mapsto \mathbb{R}$ 
pointwise if for any $x \in A$ and a given $\epsilon > 0$, we can find 
$N(x, \epsilon) \in \mathbb{N}$ such that for all $n > N(x, \epsilon)$,
$|f_n(x) - f(x)| < \epsilon$.
\end{defn}

\begin{defn}\label{s5d2}
The sequence $\{f_n\}$ converges to a function $f: A \mapsto \mathbb{R}$
uniformly if for a fixed $\epsilon > 0$, we can find $N(\epsilon) \in \mathbb{R}
$ such that for all $x \in A$ and $n > N(\epsilon)$, $|f_n(x) - f(x)| < 
\epsilon$.
\end{defn}

\begin{defn}\label{s5d3}
Let $g$ be a bounded real-valued function on $A$. Its sup-norm is defined as
\[
\lVert f\rVert_A = \sup{|g(x)|: x \in A}
\]
\end{defn}

In terms of the sup-norm, we can express the definition of uniform convergence
as: $\{f_n\}$ converges uniformly to $f$ if 
\[
\lim_{n \rightarrow \infty} \lVert f_n - f \rVert_A = 0.
\]

\begin{defn}\label{s5d4}
Given a sequence $\{f_n\}$ of functions, the series $\sum f_n$ converges if
the sequence of its partial sums $S_n(x)$ converges. The series converges
pointwise or uniformly accordingly as the sequence $\{S_n(x)\}$ converges 
pointwise or uniformly.
\end{defn}

To illustrate these ideas, we consider $A = [0, 1]$ and $f_n = x^n$. The
functions $f_n$ are bounded and continuous over their domain $A$. For all $x < 1$
$f_n(x)$ coverge pointwise to $0$ while for $x=1$ they all converge to $x=1$.
Thus, if we define $f: A \mapsto \mathbb{R}$ such that
\[
f(x) = \begin{cases}
0 & \text{ if } 0 \le x < 1 \\
1 & \text{ if } x = 1
\end{cases}
\]
then $f_n(x)$ converge to $f$. Although all $f_n$ are continuous, $f$ is not.
Thus, pointwise convergence does not preserve uniformity of the functions. It
is also evident that $\lVert f_n - f \rVert_A = 1$.

\begin{thm}[Uniform Limit Theorem]\label{s5t1}
Let $\{f_n\}$ converge to $f$ uniformly. Then if $f_n$ are continuous so is $f$.
\end{thm}
\begin{proof}
Fix an $\epsilon > 0$. Since $\{f_n\}$ converge to $f$ uniformly, we can find
$N(\epsilon) \in \mathbb{N}$ such that for all $n > N$, $|f_n(x) - f(x)| < 
\epsilon/3$ and $|f_n(x_0) - f(x_0)| < \epsilon/3$. Choose a natural number $m$
greater than $N(\epsilon)$. For this choice, $f_m$ is continuous. Therefore, we
can find $\delta > 0$ such that $|x - x_0| < \delta \Rightarrow |f(x) - f(x_0)|
< \epsilon/3$. Thus, if $|x - x_0| < \delta$,
\[
|f(x) - f(x_0)| < |f(x) - f_m(x)| + |f_m(x) - f_m(x_0)| + |f_m(x_0) - f(x_0)|
< \epsilon,
\]
\end{proof}
\begin{rem}
In this proof, we choose $\delta$ for a particular $m$. Therefore, it should
be chosen after we find $N(\epsilon)$ and select an $m$ greater than it. If all
$f_n$ have the same $\delta$ for an $\epsilon$, the sequence of functions is 
called \emph{equicontinuous}. We do not want to assume equicontinuity at this
stage.
\end{rem}

\begin{defn}[Cauchy Criterion]\label{s5d5}
Let $\{f_n\}$ be a sequence of functions. If for a given $\epsilon > 0$, we
can find $N(\epsilon) \in \mathbb{N}$ such that for all $n > m > N(\epsilon)$,
$\lVert f_n - f_m \rVert_A < \epsilon$ then $\{f_n\}$ converges uniformly.
\end{defn}

We will use the Cauchy criterion to prove,
\begin{thm}[The Weierstrass M-test]\label{s5t2}
Suppose we have a series of functions $\sum f_n$ on a set $A$. If there exist
constants $M_n$ such that $|f_n(x)| < M_n$ for all $x \in A$ and $\sum M_n$
converges then $\sum f_n$ converges uniformly and absolutely on $A$.
\end{thm}
\begin{proof}
Let $S_n(x) = f_1(x) + \cdots + f_n(x)$. Then $\sum f_n(x)$ converges if the
sequence $\{S_n(x)\}$ converges. That, $\{S_n(x)\}$ is a Cauchy sequence. Now,
if $n > m$, $|S_n(x) - S_m(x)| = |\sum_{k=m+1}^n f_k(x)| \le \sum_{k=m+1}^n
|f_k(x)| \le \sum_{k=m+1}^n M_k$. Since $\{\sum M_n\}$ converges, for a given
$\epsilon > 0$, we can find $N(\epsilon) \in \mathbb{N}$ such that for $n > m
> N(\epsilon)$, $|\sum_{k={m+1}}^n M_k| < \epsilon$. Since $\sum_{k=m+1}^n M_k 
\le |\sum_{k={m+1}}^n M_k| < \epsilon$, we have proved that for any $\epsilon 
> 0$, we can find $N(\epsilon) \in \mathbb{N}$ such that for all $n > m > 
N(\epsilon)$, $|S_n(x) - S_m(x)| < \epsilon$.
\end{proof}

\begin{thm}\label{s5t3}
Let $\{f_n\}$ be a sequence of functions Riemann-integrable on $[a, b]$. If
$f_n \rightarrow f$ uniformly then $f$ is Riemann-integrable and
\[
\int_a^b\lim_{n \rightarrow \infty} f_n(x)dx = 
\lim_{n \rightarrow \infty}\int_a^b f_n(dx).
\]
\end{thm}
\begin{proof}
Since $f_n \rightarrow f$ uniformly, for a given $\epsilon > 0$, we can find 
$N(\epsilon) \in \mathbb{N}$ such that for all $n > N(\epsilon)$, $|f_n - f| < 
\epsilon/(b - a)$ or 
\[
f - \frac{\epsilon}{b - a} < f_n < f - \frac{\epsilon}{b - a}
\]
Integrating from $a$ to $b$,
\[
\int_a^b f(x)dx - \epsilon < \int_a^b f_n(x)dx < \int_a^b f(x)dx + \epsilon
\]
for all $n > N(\epsilon)$. Therefore,
\[
\lim_{n \rightarrow \infty}\int_a^b f_n(x)dx = \int_a^b f(x)dx = 
\int_a^b \lim_{n \rightarrow \infty}f_n(x)dx.
\]
\end{proof}

\begin{rem}
Supposing the convergence was not uniform then we would have had $N$ depend on
$x$ as well as $\epsilon$. In order that $|f_n - f| < \epsilon/(b - a)$ we 
would have had to choose $N^\ast = \max(\{N(\epsilon, x) | x \in [a, b]\})$. 
The set $\mathcal{N} = \{N(\epsilon, x) | x \in [a, b]\}$ can be bounded or 
unbounded. In the former case, we can take $N^\ast=\lceil\sup\mathcal{N}\rceil$,
but then the convergence is uniform. In the latter case, no such natural number
exists and out proof cannot proceed.
\end{rem}

\begin{thm}\label{s5t5}
Let $f_n$ be a sequence of differentiable functions on $[a, b]$ such that
\begin{enumerate}
\item $f_n^\op \rightarrow g$ uniformly on $[a, b]$ and 
\item there exists at least one $x_0 \in [a, b]$ at which the sequence 
$\{f_n(x_0)\}$ of numbers converges.
\end{enumerate}
Then $f_n \rightarrow f$ uniformly and $f^\op = g$.
\end{thm}
\begin{proof}
We will first show that $\{f_n\}$ converges uniformly. Consider,
\begin{eqnarray*}
|f_n(x)-f_m(x)|&=&|f_n(x)-f_n(x_0)-f_m(x)+f_m(x_0)+f_n(x_0)-f_m(x_0)| \\
 &\le& |f_n(x) - f_n(x_0)| + |f_m(x) - f_m(x_0)| + |f_n(x_0) - f_m(x_0)|.
\end{eqnarray*}
Since the numerical sequence $\{f_n(x_0)\}$ converges, we can find $N_1 \in
\mathbb{N}$ such that for all $n, m > N_1$, $|f_n(x_0) - f_m(x_0)| < 
\epsilon/2|$. Applying the mean value theorem to $f_n - f_m$, we get
\[
|f_n(x) - f_n(x_0)| =  |f_n^\op(x^\op) - f_m^\op(x^\op)||x - x_0|
\]
for some $x^\op \in [x, x_0] (\text{ or } [x_0, x])$. Since $\{f_n^\op\}$
converge uniformly on $[a, b]$ for a given $\epsilon$, we can find $N_2 \in
\mathbb{N}$ such that for all $n, m > N_2$,
\[
|f_n^\op(x^\op) - f_m^\op(x^\op)| < \frac{\epsilon}{2|x - x_0|}
\]
We thus have,
\[
|f_n(x) - f_m(x)| < \epsilon
\]
for all $n, m > \max(N_1, N_2)$, or that $\{f_n\}$ is a Cauchy sequence. It is
uniformly convergent convergent because $N_1$ is independent of $x_0$ and $N_2$
is independent of $x^\op$ because of uniform convergence of $\{f_n^\op\}$.

Assume that $f_n \rightarrow f$. We will now show that $f^\op = g$. To that end,
consider
\begin{eqnarray*}
\Big|\frac{f(x+h) - f(x)}{h} - g(x)\Big| &=& \Big|\frac{f(x+h)-f(x)}{h} -
\frac{f_n(x+h)-f_n(x)}{h} + \\
 & & \frac{f_n(x+h)-f_n(x)}{h} - f_n^\op(x) + f_n^\op(x) - g(x)\Big| \\
  &\le& \Big|\frac{f(x+h)-f_n(x+h)}{h}\Big| + \Big|\frac{f(x)-f_n(x)}{h}\Big|\\
  & & \Big|\frac{f_n(x+h)-f_n(x)}{h} - f_n^\op(x)\Big| + |f_n^\op(x)-g(x)|
\end{eqnarray*}
Consider the term
\[
\Big|\frac{f(x+h)-f_n(x+h)}{h}\Big| + \Big|\frac{f(x)-f_n(x)}{h}\Big|.
\]
We will replace it with
\[
\Big|\frac{f_m(x+h)-f_n(x+h)}{h}\Big| + \Big|\frac{f_m(x)-f_n(x)}{h}\Big|
\]
and let $m \rightarrow \infty$. We do this because we still have not proved 
the differentiability of $f$. If $\tilde{f} = f_m - f_n$ then the above term
becomes
\[
\Big|\frac{\tilde{f}(x+h) - \tilde{f}}{h}\Big|.
\]
Since $f_n$ and $f_m$ are differentiable, so is $\tilde{f}$. Applying the mean
value theorem, this term is equal to $\tilde{f}^\op(\tilde{x})$ for some 
$\tilde{x} \in [x, x+h]$. Choose $N_1 \in \mathbb{N}$ such that for all 
$n, m > N_1$, 
\[
|f_n^\op(\tilde{x}) - f_m^\op(\tilde{x})| = |\tilde{f}^\op(
\tilde{x})| < \frac{\epsilon}{3}.
\]
In the limit $m \rightarrow \infty$, we have
\[
\Big|\frac{f(x+h)-f_n(x+h)}{h}\Big| + \Big|\frac{f(x)-f_n(x)}{h}\Big| < 
\frac{\epsilon}{3}.
\]

Since $f_n$ is differentiable in $[a, b]$, we can find $h$ such that
\begin{equation}\label{s5e1}
\Big|\frac{f_n(x+h)-f_n(x)}{h} - f_n^\op(x)\Big| < \frac{\epsilon}{3}.
\end{equation}
Finally, since $f_n^\op \rightarrow g$, we can find $N_2 \in \mathbb{N}$ such
that for all $n > N_2$, $|f_n^\op(x) - g(x)| < \epsilon/3$.

If we now choose $N = \max(N_1, N_2)$ and $h$ such that \eqref{s5e1} is 
satisfied then
\[
\Big|\frac{f(x+h) - f(x)}{h} - g(x)\Big| < \epsilon.
\]
\end{proof}

\section{Power Series}\label{s6}
\begin{defn}\label{s6d1}
A power series centred at a point $c \in \mathbb{R}$ is a function of the form
\[
\sum_{n \ge 0}a_n(x - c)^n,
\]
where $a_n \in \mathbb{R}$.
\end{defn}

\begin{defn}\label{s6d2}
The radius of convergence $R$ of a power series is a non-negative real number
such that it converges for all $|x - c| < R$ and diverges for $|x - c| > R$.
\end{defn}

\begin{thm}[Cauchy-Hadamard]\label{s6t1}
Let $\sum a_nx^n$ be a power series and let 
\[
L = \limsup \sqrt[n]{|a_n|}.
\]
Define 
\[
R = \begin{cases}
1/L & \text{ if } 0 < L < \infty \\
\infty & \text{ if } L = 0 \\
0 & \text{ if } L = \infty
\end{cases}
\]
The power series converges absolutely for $|x| < R$ and diverges for $|x| > R$.
\end{thm}
\begin{proof}
For a fixed $x_0$, let $b_n = a_n x_0^n$. Applying the root test \ref{s4p5}
to the series $\sum |b_n|$ we conclude that if $L^\op = \limsup\sqrt[n]{b_n}
= |x_0|\limsup\sqrt[n]{|a_n|}$ then the series converges if $L^\op < 1$. If
$L = \limsup\sqrt[n]{|a_n|}$ then the series converges for all $x_0$ for which
$L^\op = L|x_0| < 1$ of $|x_0| < 1/L$. 

The numerical series converges if $L^\op > 1$, that is for $|x_0| > 1/L$.
\end{proof}

\begin{thm}\label{s6t2}
If a power series $\sum a_nx^n$ has a radius of convergence $R > 0$, then for
$0 < r < R$, then the series converges uniformly in $[-r, r]$.
\end{thm}
\begin{proof}
Choose $r$ such that $0 < r < R$ and define $M_n = |a_n|r^n$. $\sum M_n$ 
converges. Since $a_nx^n \le |a_n|x^n < |a_n|r^n = M_n$ by Weierstrass M-test, 
$\sum a_nx^n$ converges uniformly for any $x \in [-r, r]$.
\end{proof}
\end{document}

